
\Slide{1}{连续性在本课程中的作用}
函数连续描述输入小变动不会造成输出的大跳跃。本讲从数列判据开始，逐步得到紧集上的极值与连通集上的介值，再研究只有一个方向受控制的半连续性。

\Slide{2}{连续性在本课程中的作用}
本页继续“连续性在本课程中的作用”。完整定义、条件和论证见本讲原页第 1 页下方。

\Slide{3}{序列连续与复合、坐标}
设 $f:S\to T$，在 $x$ 连续指每个 $x_n\in S,x_n\to x$ 都满足 $f(x_n)\to f(x)$。复合连续：先由 $f$ 连续得 $f(x_n)\to f(x)$，再由 $g$ 连续取像。有限维中 $\norm{z_n-z}\to0$ 等价于每个坐标收敛：一边用 $|z_{ni}-z_i|\le\norm{z_n-z}$，另一边用有限和 $\norm{z_n-z}^2=\sum_i|z_{ni}-z_i|^2$。所以向量函数连续当且仅各坐标函数连续。

\Slide{4}{序列连续与复合、坐标}
本页继续“序列连续与复合、坐标”。完整定义、条件和论证见本讲原页第 3 页下方。

\Slide{5}{序列连续与复合、坐标}
本页继续“序列连续与复合、坐标”。完整定义、条件和论证见本讲原页第 3 页下方。

\Slide{6}{序列连续与复合、坐标}
本页继续“序列连续与复合、坐标”。完整定义、条件和论证见本讲原页第 3 页下方。

\Slide{7}{序列连续与复合、坐标}
本页继续“序列连续与复合、坐标”。完整定义、条件和论证见本讲原页第 3 页下方。

\Slide{8}{连续性的五个等价条件}
连续的 $\eps$--$\delta$ 形式为：对每个 $\eps>0$ 存在 $\delta>0$，所有定义域内距 $x$ 小于 $\delta$ 的点，其像距 $f(x)$ 小于 $\eps$。否定它可逐个取 $\delta=1/n$ 构造 $x_n\to x$ 而像距始终至少 $\eps_0$，所以与数列定义等价。连续推出开集逆像开：$f(x)\in G$ 先取像球包含于 $G$，再用连续性拉回一个输入球。开逆像与闭逆像用补集对偶。闭逆像推出 $f(\overline A)\subseteq\overline{f(A)}$：后者的逆像是含 $A$ 的闭集，故含闭包。反向若有闭包包含性质，取 $A=f^{-1}(F)$ 且 $F$ 闭，则 $f(\overline A)\subseteq F$，故 $\overline A\subseteq A$，逆像闭。不要把像闭包上的横线漏掉。

\Slide{9}{连续性的五个等价条件}
本页继续“连续性的五个等价条件”。完整定义、条件和论证见本讲原页第 8 页下方。

\Slide{10}{连续性的五个等价条件}
本页继续“连续性的五个等价条件”。完整定义、条件和论证见本讲原页第 8 页下方。

\Slide{11}{连续性的五个等价条件}
本页继续“连续性的五个等价条件”。完整定义、条件和论证见本讲原页第 8 页下方。

\Slide{12}{实值连续函数的水平集}
实值函数连续当且仅所有 $\{f\le a\}$ 与 $\{f\ge a\}$ 闭。必要性是闭集逆像；充分性先取补得 $\{f>a\},\{f<b\}$ 开，再相交得 $f^{-1}((a,b))$ 开，任意实开集为开区间并，所以逆像开。实开集的极大区间分解见106第53页：不交区间各含不同有理数，因此可数。

\Slide{13}{实值连续函数的水平集}
本页继续“实值连续函数的水平集”。完整定义、条件和论证见本讲原页第 12 页下方。

\Slide{14}{实值连续函数的水平集}
本页继续“实值连续函数的水平集”。完整定义、条件和论证见本讲原页第 12 页下方。

\Slide{15}{紧集的连续像与极值}
设 $K$ 非空紧，$\{U_i\}$ 覆盖 $f(K)$，则 $\{f^{-1}(U_i)\}$ 为 $K$ 的开覆盖。取有限子覆盖再取像得 $f(K)$ 的有限子覆盖，故像紧。实值时像非空闭有界，令 $M=\sup f(K)$，选像中数 $y_n>M-1/n$，闭性给 $M\in f(K)$，故有 $x_M\in K$ 达到最大。最小同理。$M$ 是值，$x_M$ 是达到该值的点；最优点可以不唯一。

\Slide{16}{紧集的连续像与极值}
本页继续“紧集的连续像与极值”。完整定义、条件和论证见本讲原页第 15 页下方。

\Slide{17}{紧集的连续像与极值}
本页继续“紧集的连续像与极值”。完整定义、条件和论证见本讲原页第 15 页下方。

\Slide{18}{连通像与介值}
若 $f(K)=A\cup B$ 是两个非空不交相对开集的分离，则逆像为 $K$ 的非空不交相对开分离，矛盾。实数连通集是区间，所以连续函数的两个函数值之间的每个数也是某个函数值。在 $[a,b]$ 上，若 $f(a)f(b)<0$ 就存在零点。仅能得存在，不保证唯一；唯一通常还需严格单调等条件。

\Slide{19}{连通像与介值}
本页继续“连通像与介值”。完整定义、条件和论证见本讲原页第 18 页下方。

\Slide{20}{邻域语言与量词}
邻域包含一个以该点为中心的开球。连续要求每个输出邻域的逆像都是输入点的邻域，不要求逆像本身是球。量词顺序为 $\forall\eps\ \exists\delta\ \forall z$，允许 $\delta$ 依赖 $x$ 与 $\eps$，不能在看到 $z$ 后再选它。例 $f(t)=t^2$ 在 $x$ 附近取 $\delta=\min(1,\eps/(2|x|+1))$，因 $|z^2-x^2|\le|z-x|(2|x|+1)$。

\Slide{21}{邻域语言与量词}
本页继续“邻域语言与量词”。完整定义、条件和论证见本讲原页第 20 页下方。

\Slide{22}{邻域语言与量词}
本页继续“邻域语言与量词”。完整定义、条件和论证见本讲原页第 20 页下方。

\Slide{23}{实值半连续的定义与等价判据}
下半连续指 $f(x)\le\liminf f(x_n)$，允许向下的低点；上半连续指 $f(x)\ge\limsup f(x_n)$，允许向上的高点。下半连续等价于 $\{f>a\}$ 开，也等价于 $\{f\le a\}$ 闭：若 $a<f(x)$，定义给邻域中 $f>a$；反向对任意 $a<f(x)$ 用该开集取邻域。上半连续用 $-f$ 转化，得到 $\{f<a\}$ 开及 $\{f\ge a\}$ 闭。两者同时成立即连续。例 $f(0)=0,f(x)=1$ 对 $x\ne0$ 下半连续但不连续，其负函数上半连续。

\Slide{24}{实值半连续的定义与等价判据}
本页继续“实值半连续的定义与等价判据”。完整定义、条件和论证见本讲原页第 23 页下方。

\Slide{25}{实值半连续的定义与等价判据}
本页继续“实值半连续的定义与等价判据”。完整定义、条件和论证见本讲原页第 23 页下方。

\Slide{26}{实值半连续的定义与等价判据}
本页继续“实值半连续的定义与等价判据”。完整定义、条件和论证见本讲原页第 23 页下方。

\Slide{27}{实值半连续的定义与等价判据}
本页继续“实值半连续的定义与等价判据”。完整定义、条件和论证见本讲原页第 23 页下方。

\Slide{28}{实值半连续的定义与等价判据}
本页继续“实值半连续的定义与等价判据”。完整定义、条件和论证见本讲原页第 23 页下方。

\Slide{29}{实值半连续的定义与等价判据}
本页继续“实值半连续的定义与等价判据”。完整定义、条件和论证见本讲原页第 23 页下方。

\Slide{30}{半连续函数的极值证明}
非空紧 $K$ 上有限实值下半连续函数先必有下界：若无，取 $f(x_n)<-n$，紧性给子列趋于 $x\in K$，下半连续将推出有限 $f(x)\le-\infty$，矛盾。设有限 $m=\inf_K f$，选 $f(x_n)<m+1/n$，再取收敛子列，得到 $f(x)\le\liminf f(x_n)=m$，而下界给反向，故取到最小值。对 $-f$ 应用此结论得上半连续函数有最大值。原页35趋近的值应为 $M$ 而非 $m$。不能只从紧集直接说任意函数有界。

\Slide{31}{半连续函数的极值证明}
本页继续“半连续函数的极值证明”。完整定义、条件和论证见本讲原页第 30 页下方。

\Slide{32}{半连续函数的极值证明}
本页继续“半连续函数的极值证明”。完整定义、条件和论证见本讲原页第 30 页下方。

\Slide{33}{半连续函数的极值证明}
本页继续“半连续函数的极值证明”。完整定义、条件和论证见本讲原页第 30 页下方。

\Slide{34}{半连续函数的极值证明}
本页继续“半连续函数的极值证明”。完整定义、条件和论证见本讲原页第 30 页下方。

\Slide{35}{半连续函数的极值证明}
本页继续“半连续函数的极值证明”。完整定义、条件和论证见本讲原页第 30 页下方。

\Slide{36}{上确界、下确界与有限运算}
若 $g=\sup_i f_i$ 且各 $f_i$ 下半连续，则 $\{g>a\}=\bigcup_i\{f_i>a\}$ 开，故 $g$ 下半连续。上半连续函数的任意下确界对偶。有限下半连续函数的最小值满足 $\{\min_i f_i>a\}=\bigcap_i\{f_i>a\}$，有限交开；上半连续函数有限最大值对偶。无限交不保证开，故不能删掉“有限”。若确界可能无穷，应注明把函数视为扩展实值。

\Slide{37}{上确界、下确界与有限运算}
本页继续“上确界、下确界与有限运算”。完整定义、条件和论证见本讲原页第 36 页下方。

\Slide{38}{上确界、下确界与有限运算}
本页继续“上确界、下确界与有限运算”。完整定义、条件和论证见本讲原页第 36 页下方。

\Slide{39}{上确界、下确界与有限运算}
本页继续“上确界、下确界与有限运算”。完整定义、条件和论证见本讲原页第 36 页下方。

\Slide{40}{半连续函数的加法与正乘法}
下半连续 $f,g$ 的和下半连续：给定 $\eps$，分别取邻域使 $f(z)>f(x)-\eps/2$ 和 $g(z)>g(x)-\eps/2$，在交邻域相加。对上半连续用负号对偶。正函数的乘积可用单调连续的 $\log$：$\log f,\log g$ 保持相应半连续性，和半连续，再用单调连续的指数还原。也可直接给乘积误差界。正性不能无故删掉，因为乘负数会反转半连续的方向。

\Slide{41}{半连续函数的加法与正乘法}
本页继续“半连续函数的加法与正乘法”。完整定义、条件和论证见本讲原页第 40 页下方。

\Slide{42}{半连续函数的加法与正乘法}
本页继续“半连续函数的加法与正乘法”。完整定义、条件和论证见本讲原页第 40 页下方。

\Slide{43}{半连续函数的加法与正乘法}
本页继续“半连续函数的加法与正乘法”。完整定义、条件和论证见本讲原页第 40 页下方。
